\section{Related Literature}\label{sec:literature}

This thesis sits where four literatures meet: the potential-outcomes approach to dynamic causal effects in
macroeconomics, the identification of monetary policy shocks, the use of the Federal Reserve's text as data, and
the econometrics of weighting estimators when overlap is limited or fails. Each has moved substantially since the
first version of this thesis was written in 2020.

\subsection{Dynamic causal effects and potential outcomes}

\citet{angrist2011} and \citet{angrist2018} recast the effect of monetary policy as a policy-evaluation problem:
each FOMC decision is a treatment, a propensity score for the Committee's choice is estimated by ordered probit,
and impulse responses are inverse-probability-weighted averages of future outcomes. Their approach is robust to
nonlinearity and asymmetry, and it requires selection on observables and overlap. \citet{rambachan2025} give the
general framework: a ``direct potential outcome system'' in which impulse responses and local projections have a
nonparametric causal interpretation under conditions they make precise. This thesis adopts their notation for
potential outcomes indexed by the policy action and the horizon.

A closely related line asks what the standard linear tools estimate when the world is not linear.
\citet{plagborg2021} show that local projections and VARs estimate the same population impulse responses, so the
choice between them is about finite-sample bias and variance, not about the estimand. \citet{kolesar2025} show that
VARs and linear local projections on an observed shock or proxy identify \emph{weighted averages} of causal effects
however nonlinear the data-generating process, whereas methods that rely on heteroskedasticity or non-Gaussianity
for identification can fail. The IPW estimator used here targets an unweighted average effect of a discrete policy
action instead. The price is that it needs a correctly specified propensity score and, as this thesis emphasises,
overlap. \citet{montielolea2021} show that lag-augmented local projections give valid inference even with persistent
data and at long horizons. Their result is a useful benchmark for the inference problem this thesis faces, in which
outcomes at nearby meetings overlap in time.

\subsection{Identifying monetary policy shocks and the Federal Reserve's information}

The central difficulty in estimating the effect of monetary policy is that the Federal Reserve reacts to what it
expects the economy to do. \citet{romer2004} addressed this by purging the intended funds-rate change of its
response to the Greenbook forecasts. \citet{romer2023} revisit thirty-five years of narrative evidence and argue
that carefully identified shocks have large effects on output, unemployment and inflation in the expected
directions. High-frequency identification uses the change in market rates in a narrow window around FOMC
announcements instead. \citet{nakamura2018} argue that such surprises also reveal the Fed's private information
about the economy (an ``information effect''), and \citet{jarocinski2020} and \citet{mirandaagrippino2021}
separate pure policy shocks from information shocks using the joint response of stock prices or private forecasts.
\citet{bauer2023aer,bauer2023nber} offer an alternative account: much of the apparent information effect reflects
the Fed responding to economic news that markets had not yet priced, which makes raw surprises predictable and
calls for orthogonalising them with respect to prior news.

Most closely related to this thesis, \citet{aruoba2025} use natural language processing on the staff documents
prepared before each FOMC meeting to measure the Committee's information set. They find that the text contains
information about the economy that the Greenbook's numerical forecasts do not, and that shocks identified as
residuals from a model that uses it produce responses consistent with standard theory. Their conclusion speaks
directly to the empirical results of Section~\ref{sec:results}: conditioning on numerical forecasts alone may leave
enough of the Fed's information out of the propensity score to bias selection-on-observables estimates.

\subsection{The Federal Reserve's text as data}

\citet{hansen2018} brought computational linguistics to the FOMC transcripts, and the use of the Fed's internal
documents has since grown quickly. \citet{ke2022} develop robust topic-model estimators and apply them to the
transcripts of the Greenspan era, the period studied here. Two recent papers use the Bluebook alternatives
themselves. \citet{doh2026} fine-tune a language model on the alternative policy statements drafted by Board staff
to measure where the adopted statement sits on the dovish--hawkish spectrum, and \citet{laarits2025} use the
alternatives circulated before each meeting to study how the Committee aggregates members' differing economic
models. This thesis uses the same documents for a different purpose: not to measure the tone of communication, but
to recover the set of actions the Committee considered, which determines where the propensity score is zero.

\subsection{Overlap, weighting and partial identification}

When propensity scores approach zero, IPW estimators become unstable and standard inference can fail.
\citet{khan2010} show that identification is then irregular and convergence slower than $\sqrt{n}$;
\citet{crump2009} propose trimming to a subpopulation with good overlap; \citet{ma2020} and \citet{heiler2021}
develop inference that remains valid with small or many extreme scores; and \citet{damour2021} show that strict
overlap becomes restrictive when there are many covariates. All of these treat limited overlap as a feature of the
estimated score. The distinguishing feature of the present setting is that overlap fails \emph{exactly}, and the
region where it fails is recorded in the data.

For that case, partial identification is the natural tool. Worst-case bounds under a support restriction go back to
\citet{manski1990,manski2003}, and \citet{imbens2004} give confidence intervals for the partially identified
parameter. \citet{susmann2025} develop bounds on the average treatment effect in the cross-sectional setting when
outcome overlap fails for a subpopulation, with width equal to the outcome range times the probability of that
subpopulation, and propose efficient estimators and uniformly valid inference. Proposition~\ref{prop:bounds}
is the analogue for a multi-valued treatment in a time series, where the non-overlap region is observed rather than
estimated and some menus contain only one of the two actions being compared. Finally, a large literature relaxes
unconfoundedness itself; \citet{masten2018} bound treatment effects when selection on unobservables is limited by a
sensitivity parameter, an approach that complements the $\lambda$-relaxation of Section~\ref{sec:framework}, which
relaxes effect homogeneity across menus instead.
